\documentclass[tikz,border=12pt]{standalone}
% Shared visual language for standalone EMT block diagrams.
\usepackage{amsmath}
\usetikzlibrary{arrows.meta,backgrounds,calc,fit,positioning}

\tikzset{
    emt diagram/.style = {>={Stealth[length=2.4mm]}},
    line/.style        = {semithick, rounded corners=6pt},
    block/.style       = {draw, semithick, fill=white,
                          minimum height=1cm, minimum width=2.3cm},
    hblock/.style      = {block, minimum width=2.24cm, minimum height=1.3cm},
    vblock/.style      = {block, minimum width=1.3cm, minimum height=2.24cm},
    sum/.style         = {draw, semithick, fill=white, circle,
                          inner sep=2.6pt},
    signal/.style      = {circle, fill, inner sep=1.4pt,
                          label distance=3pt},
    sig/.style         = {font=\small},
    pm/.style          = {font=\scriptsize},
    wire jump/.style   = {line,
                          preaction={draw=white, line width=3pt,
                                     line cap=round}},
    enclosure/.style   = {draw, dashed, semithick, rounded corners=4pt},
    operator enclosure/.style = {enclosure, inner xsep=7mm, inner ysep=6mm}
}


% TUNABLES: gain width and spacing of the two lower branches.
\def\gainwidth{1.35cm}
\def\integralrow{-1.9}
\def\trackingrow{-3.8}

\begin{document}
\begin{tikzpicture}[emt diagram]

% Current reference and proportional path.
\node[block, minimum width=2.9cm, minimum height=1.3cm]
    (Ilim) at (-0.4,0)
    {$\dfrac{\mathbf{i}^{\mathrm{ref}}}
        {\sqrt{\mathcal{L}_i(\mathbf{i}^{\mathrm{ref}})}}$};
\node[signal] (ilim) at (1.65,0) {};
\node[sum] (error) at (3,0) {$\Sigma$};
\coordinate (esplit) at (4.2,0);
\node[block, minimum width=\gainwidth] (Kp) at (6,0) {$K_P$};
\node[sum] (command) at (10.2,0) {$\Sigma$};
\draw[line] (Ilim.east) -- (ilim);
\draw[->, line] (ilim) -- (error.west);
\node[pm, above left=0pt and 1pt of error.west] {$+$};
\draw[line] (error.east) -- (esplit);
\draw[->, line] (esplit) -- node[sig, above] {$\mathbf{e}$} (Kp.west);
\draw[->, line] (Kp.east) -- (command.west);
\node[pm, above left=0pt and 1pt of command.west] {$+$};

% Integral path, with anti-windup added before integration.
\node[block, minimum width=\gainwidth] (Ki)
    at (4.2,\integralrow) {$K_I$};
\node[sum] (rate) at (6,\integralrow) {$\Sigma$};
\node[block, minimum width=\gainwidth] (integrator)
    at (8.1,\integralrow) {$\dfrac{\mathbf{I}}{s}$};
\node[signal, label={[sig]above:$\boldsymbol{\xi}$}] (xi)
    at (9.35,\integralrow) {};
\draw[->, line] (esplit) -- (Ki.north);
\draw[->, line] (Ki.east) -- (rate.west);
\node[pm, above left=0pt and 1pt of rate.west] {$+$};
\draw[->, line] (rate.east) -- (integrator.west);
\draw[line] (integrator.east) -- (xi);
\draw[->, line] (xi) -| (command.south);
\node[pm, below right=1pt and 0pt of command.south] {$+$};

% Capacitor-voltage feedforward and synchronous-frame decoupling.
\node[block, minimum width=4cm, minimum height=1.7cm]
    (feedforward) at (10.2,3)
    {$\mathbf{v}+\mathcal{J}\omega L\mathbf{i}$};
\draw[->, line] (feedforward.south)
    -- node[sig, right] {$\mathbf{b}$} (command.north);
\node[pm, above right=1pt and 0pt of command.north] {$+$};

% DC-dependent voltage-command limit.
\node[signal, label={[sig]above:$\mathbf{z}$}] (z) at (11.35,0) {};
\node[block, minimum width=3.2cm, minimum height=1.3cm]
    (Ulim) at (13.6,0)
    {$\dfrac{v_{\mathrm{dc}}\mathbf{z}}
        {\sqrt{\mathcal{L}_u(v_{\mathrm{dc}},\mathbf{z})}}$};
\node[signal, label={[sig]above:$\mathbf{u}$}] (u) at (16.3,0) {};
\draw[line] (command.east) -- (z);
\draw[->, line] (z) -- (Ulim.west);
\draw[line] (Ulim.east) -- (u);
\draw[->, line] (u) -- (18,0);

% Tracking anti-windup: limited minus unlimited voltage command.
\node[sum] (tracking) at (11.35,\trackingrow) {$\Sigma$};
\node[block, minimum width=\gainwidth] (Kaw)
    at (8.1,\trackingrow) {$K_{\mathrm{aw}}$};
\coordinate (uloop) at (16.3,-4.9);
\draw[->, line] (z) -- (tracking.north);
\node[pm, above left=1pt and 0pt of tracking.north] {$-$};
\draw[->, line] (u) -- (uloop) -| (tracking.south);
\node[pm, below left=1pt and 0pt of tracking.south] {$+$};
\draw[->, line] (tracking.west) -- (Kaw.east);
\draw[->, line] (Kaw.west) -| (rate.south);
\node[pm, below left=1pt and 0pt of rate.south] {$+$};

% External inputs share the same frame and signal ownership conventions.
\node[signal, label={[sig]above:$\mathbf{i}^{\mathrm{ref}}$}]
    (iref) at (-3.8,0) {};
\node[signal, label={[sig]above:$\mathbf{i}$}] (i) at (-3.8,2.5) {};
\node[signal, label={[sig]above:$\omega$}] (omega) at (-3.8,3) {};
\node[signal, label={[sig]above:$\mathbf{v}$}] (v) at (-3.8,3.5) {};
\node[signal, label={[sig]right:$v_{\mathrm{dc}}$}]
    (vdc) at (13.6,5.4) {};
\coordinate (isplit) at (error |- i);
\draw[->, line] ($(iref)+(-0.8,0)$) -- (Ilim.west);
\draw[->, line] ($(i)+(-0.8,0)$) -- (feedforward.west |- i);
\draw[->, line] (isplit) -- (error.north);
\node[pm, above left=1pt and 0pt of error.north] {$-$};
\draw[->, line] ($(omega)+(-0.8,0)$) -- (feedforward.west |- omega);
\draw[->, line] ($(v)+(-0.8,0)$) -- (feedforward.west |- v);
\draw[->, line] ($(vdc)+(0,0.8)$) -- (Ulim.north);

% The outer voltage loop can track the limited current reference.
\draw[->, line] (ilim) -- (1.65,-6.2);
\node[sig, right] at (1.65,-5.9) {$\mathbf{i}^{\mathrm{lim}}$};

% Enclose the model-owned states, variables, and feedback paths.
\begin{scope}[on background layer]
    \node[operator enclosure,
        fit=(Ilim)(ilim)(error)(Kp)(Ki)(rate)(integrator)(xi)
            (feedforward)(command)(z)(Ulim)(u)(tracking)(Kaw)(uloop)] {};
\end{scope}

\end{tikzpicture}
\end{document}
